\documentclass[11pt]{article}
\usepackage{amsmath,amssymb,amsthm}
\usepackage[margin=1in]{geometry}
\newtheorem{theorem}{Theorem}
\newtheorem{lemma}{Lemma}
\newtheorem{proposition}{Proposition}
\newtheorem{corollary}{Corollary}
\title{V76: Top-Tree Transfer for Support-Boundary Decompositions}
\author{NC0\_k-Avoid Laboratory}
\date{August 2026}
\begin{document}
\maketitle

\paragraph{Setting.}
For a set $M$ of gate supports, define
\[
 \lambda(S)=\left|\left(\bigcup_{i\in S}\operatorname{supp}(i)\right)
 \cap\left(\bigcup_{i\notin S}\operatorname{supp}(i)\right)\right|.
\]
A supplied branch decomposition is an unrooted subcubic tree with gate-labelled
leaves and width at most $b$ under $\lambda$.

\paragraph{Prior-art ingredient.}
A labelled top tree is a binary hierarchy of connected clusters with at most
two boundary vertices. Standard top-tree constructions have logarithmic
height on an underlying labelled tree of linear size.

\begin{lemma}[Four-cut cover lemma]
Let $C$ be a labelled top-tree cluster of a subcubic supplied branch tree, and
let $S(C)$ be the gate labels included in $C$. There are at most four original
branch-tree edges whose middle sets cover the support boundary of $S(C)$.
Consequently,
\[
 \lambda(S(C))\le 4b.
\]
\end{lemma}

\begin{proof}
The cluster has at most two boundary vertices. At each such vertex, at most
two incident tree edges or label crossings separate cluster atoms from their
complement. A gate-label crossing at an original leaf is represented by that
leaf's unique incident branch edge. Hence at most four original edges suffice.
For any input variable occurring both in $S(C)$ and outside $S(C)$, the unique
path between corresponding gate leaves crosses the labelled-cluster boundary,
and therefore one selected original edge. The variable lies in that edge's
middle set.
\end{proof}

\begin{theorem}[Top-tree support-boundary transfer]
From a supplied width-$b$ subcubic gate branch decomposition, one can construct
a rooted binary gate decomposition $T'$ satisfying
\[
 \operatorname{width}(T')\le 4b,
 \qquad
 \operatorname{height}(T')=O(\log m),
 \qquad
 \operatorname{EPL}(T')=O(m\log m).
\]
\end{theorem}

\begin{proof}
Attach each gate as a unique label to its branch-tree leaf and take a standard
binary labelled top tree. Delete base leaves carrying no gate label and
suppress unary nodes. Height does not increase, every retained node remains a
labelled cluster, and the retained leaves are exactly the gates. Apply the
four-cut cover lemma at every retained node. The external-path-length bound
follows from logarithmic height and $m$ leaves.
\end{proof}

\begin{corollary}[V75 consequence]
Rebuilding the V75 monotone residual circuit on $T'$ gives incremental prefix
avoidance time
\[
 O\!\left(m\log m\,A(4b)^2\operatorname{poly}(n,m)\right).
\]
\end{corollary}

\begin{proposition}[Perfect-height tradeoff witness]
For the seven rank-two supports
\[
 \{0\},\{1\},\{0,2\},\{0,3\},\{1,2\},\{1,3\},\{2,3\},
\]
the exact nondominated $(\mathrm{width},\mathrm{height},\mathrm{EPL})$ frontier
is
\[
 (2,4,21),\qquad (3,3,20).
\]
\end{proposition}

\paragraph{Scope.}
The top-tree height theorem is prior art. No optimality of the factor four,
width-preserving logarithmic impossibility, unrestricted range-avoidance
algorithm, standard-model lower bound, novelty, peer review, or P-versus-NP
resolution is claimed. A previous centroid-based factor-two proof attempt was
rejected and is not a theorem of this module.

\end{document}
